\section{Positive dimension for the sporadic family}
\label{familypos:section}

The absence of a positive finite integer witness does not imply that
the corresponding $3$-adic intersection has dimension zero. We now
determine this distinction for every multiplier in
Theorem~\ref{thm:family}. Put
\[
 M_k=621\cdot3^{4k}-20,\qquad
 (m_k,L_k)=
 \begin{cases}(3,12),&k=0,\\(4k+3,12k+6),&k\ge1.\end{cases}
\]
We use $\mathcal H=\{M\in\N:\dim_H C(1,M)>0\}$ for the
multiplier set counted by $\MH(X)$; $\mathcal V$ remains the set
of positive finite-integer quotients counted by $\MC(X)$.

\begin{theorem}[Dimension of an explicit subsystem]
\label{familypos:main}
For every $k\ge0$, let $d_k>0$ be the unique solution of
\begin{equation}\label{familypos:eq:root}
             3^{-m_kd_k}+3^{-L_kd_k}=1.
\end{equation}
There is a compact subset of $C(1,M_k)$ of Hausdorff dimension
exactly $d_k$. Consequently
\[
             \dim_H C(1,M_k)\ge d_k>0,
             \qquad M_k\in\mathcal H\setminus\mathcal V.
\]
If $\vartheta>1$ is the positive root of
$\vartheta^4-\vartheta^3-1=0$, then
\begin{equation}\label{familypos:eq:theta}
 d_k\ge\frac{\log_3\vartheta}{m_k},
\end{equation}
with equality for $k=0$ and strict inequality for $k\ge1$.
\end{theorem}

The exact value in~\eqref{familypos:eq:root} concerns the constructed
subsystem, not all of $C(1,M_k)$. The conclusion that these multipliers
are not finite-integer quotients uses Theorem~\ref{thm:family}; the
positive-dimension construction below is independent of its proof.
The distinction between positive dimension and finite witnesses is
already present in Lagarias's work~\cite{Lagarias}, including the
example $M=52$. Our assertion is the explicit construction and
dimension bound for this particular sporadic family proposed
in~\cite[\S1.4]{BMRS}.

\subsection{Words and carry paths}

In this section a word $w=w_0\cdots w_{\ell-1}$ is written
\emph{least significant digit first}, with
\[
 |w|=\ell,\qquad [w]=\sum_{j=0}^{\ell-1}w_j3^j.
\]
Juxtaposition denotes concatenation, so $[uv]=[u]+3^{|u|}[v]$.
All displayed input and output words have digits in $\{0,1\}$.
A carry transition for multiplication by $M\in\N$ is
\begin{equation}\label{familypos:eq:carry}
             Ma_j+c_j=b_j+3c_{j+1}.
\end{equation}

\begin{lemma}[An integer identity certifies a path]
\label{familypos:lem:path}
Let $c,d\ge0$ be integers, and let $w,v$ be binary words of
the same length $\ell$. If
\begin{equation}\label{familypos:eq:path}
                 M[w]+c=[v]+3^\ell d,
\end{equation}
then digitwise multiplication gives an admissible path from carry
$c$ to carry $d$, with input $w$ and output $v$.
\end{lemma}

\begin{proof}
Reduction modulo $3$ identifies the first output digit with $v_0$.
Since $Mw_0+c\ge0$, the quotient after removing this residue is a
nonnegative integer. Subtract the first-digit equality from
\eqref{familypos:eq:path} and divide by $3$. Repetition proves the
claim at every digit, including the prescribed final carry.
\end{proof}

Suppose a fixed admissible prefix from carry $0$ reaches $c$, and
two words $A,B$ return $c$ to itself. Any infinite concatenation of
these two words, appended to the prefix, is an admissible input
$x\in\Sigma$. Its output digits also lie in $\{0,1\}$, and they
represent $Mx$: telescoping~\eqref{familypos:eq:carry} through the
first $n$ digits leaves the error $3^nc_n$, which tends to zero in
$\Z_3$ because $c_n$ is an integer. Thus the whole construction
lies in $C(1,M)$. It is essential here that the return carry is
reachable from $0$.

\subsection{The initial multiplier}

For $k=0$ we have $M_0=601$. Take $c=69$ and the following paths:
\[
\begin{array}{c|c|c|c}
 &\text{input}&\text{output}&\text{carries}\\\hline
 \text{prefix}&1110010010&1011001101&0\longrightarrow69\\
 A&010&001&69\longrightarrow69\\
 B&101110010010&110101101101&69\longrightarrow69
\end{array}
\]
The respective input values are $6817,3,61354$, and the output
values are $22636,9,204394$. The three exact identities are
\begin{align*}
 601\cdot6817&=22636+69\cdot3^{10},\\
 601\cdot3+69&=9+69\cdot3^3,\\
 601\cdot61354+69&=204394+69\cdot3^{12}.
\end{align*}
Lemma~\ref{familypos:lem:path} therefore supplies a reachable
carry with two return words of lengths $3$ and $12$, starting
with different digits.

\subsection{A reachable carry for every positive parameter}

Fix $k\ge1$ and write
\[
 t=81^k,\quad m=4k+3,\quad L=12k+6,\quad
 M=621t-20=23\cdot3^m-20,\quad c=69t.
\]
The prefix
\begin{equation}\label{familypos:eq:prefix}
        P=(1100)^k\,111\,0^{4k}\,1\,0^{4k+2}
\end{equation}
has length $L$. Its value is
\[
 p=[P]=\frac{t-1}{20}+13t+27t^2
       =\frac{540t^2+261t-1}{20}.
\]
Using $3^L=729t^3$, direct multiplication gives
\begin{equation}\label{familypos:eq:prefix-output}
 q:=Mp-23\cdot3^L
      =\frac{151281t^2-5841t+20}{20}.
\end{equation}
We next verify that $q$ is the value of a binary ternary word of
length $L$, allowing zeros at its high end. For $k=1$, its base-$81$
digits in least-significant-first order are
\[
                         (1,36,27,12,1).
\]
For $k\ge2$ they are
\[
 \bigl(1,\underbrace{0,\ldots,0}_{k-1},36,0,
       \underbrace{4,\ldots,4}_{k-2},31,12,1\bigr).
\]
Indeed, the latter digit list is the identity
\begin{align*}
 q={}&1+36\cdot81^k
       +4\sum_{j=k+2}^{2k-1}81^j\\
    &+31\cdot81^{2k}+12\cdot81^{2k+1}+81^{2k+2},
\end{align*}
which follows from the geometric sum formula and
\eqref{familypos:eq:prefix-output}; an empty sum is zero.
Every listed base-$81$ digit belongs to
$\{0,1,4,12,27,31,36\}$ and hence has four ternary digits
in $\{0,1\}$ after padding. The highest nonzero base-$81$ digit
is $1$ in position $2k+2$, so the ternary length is at most
$8k+9\le L$. Thus $P$ reaches carry $23$.

The input $10$ with output $01$ then reaches $c$, since
\begin{equation}\label{familypos:eq:entry}
                  M+23=3+9c.
\end{equation}
Consequently the binary prefix $P10$ reaches $c=69t$ from the
initial carry $0$.

\subsection{Two return cycles}

The short cycle is
\[
 A=0^{m-2}10,\qquad |A|=m,\qquad [A]=3^{m-2}=3t.
\]
Its output is $0^{m-1}1$, because
\begin{equation}\label{familypos:eq:short-cycle}
                    M[A]+c=3^{m-1}+3^mc.
\end{equation}
For the long cycle, put
\begin{equation}\label{familypos:eq:long-word}
 B=(1100)^{k-1}\,11010110\,(0110)^{k-1}\,01\,
                              0^{4k+2}\,10.
\end{equation}
The lengths of its six successive pieces sum to $L$. To compute
its value, write $r=81^{k-1}=t/81$ and use concatenation:
\begin{align*}
 b=[B]
   &=\frac{r-1}{20}+1003r
      +\frac{19683}{20}r(r-1)+3t^2+81t^3\\
   &=81t^3+\frac{63}{20}t^2+\frac7{30}t-\frac1{20}.
\end{align*}
Here $[11010110]=1003$ and $[0110]=12$; the two geometric
sums account for the repeated four-digit pieces.
The required output value is
\begin{equation}\label{familypos:eq:long-output}
 z:=Mb+c-3^Lc
       =\frac{20169t^3+4914t^2+1997t+60}{60}.
\end{equation}
For every $k\ge1$, the geometric sum formula gives the exact
base-$81$ expression
\begin{align*}
 z={}&1+27\cdot81^{k-1}+37\cdot81^k
      +4\sum_{j=k+1}^{2k-1}81^j\\
    &+13\cdot81^{2k}+13\cdot81^{2k+1}
      +12\sum_{j=2k+2}^{3k}81^j+4\cdot81^{3k+1}.
\end{align*}
For $k\ge2$ the nonoverlapping terms give the digit list
\[
 \bigl(1,\underbrace{0,\ldots,0}_{k-2},27,37,
       \underbrace{4,\ldots,4}_{k-1},13,13,
       \underbrace{12,\ldots,12}_{k-1},4\bigr).
\]
For $k=1$ only the first two terms overlap, combining to $28$;
the list is $(28,37,13,13,4)$. All the digits in either list
belong to $\{0,1,4,12,13,27,28,37\}$, each of which has only
$0,1$ as ternary digits. The highest digit $4=(11)_3$ is in
base-$81$ position $3k+1$ and occupies ternary positions
$L-2,L-1$. Therefore $z<3^L$ and $z$ is a valid length-$L$
output word.

Lemma~\ref{familypos:lem:path}, applied to
\eqref{familypos:eq:short-cycle} and
\eqref{familypos:eq:long-output}, proves that both $A$ and $B$
return to $c$. Their first digits are respectively $0$ and $1$.
These identities and digit patterns apply to all $k\ge1$; no
extrapolation from finitely many carry graphs is used.

\subsection{Dimension of the return-word construction}

We give the dimension argument explicitly, including the differing
word lengths. For either construction above, let $K\subset\Z_3$
be the set of all infinite concatenations of $A$ and $B$ without
the fixed entry prefix. This is compact, and
\[
        K=([A]+3^mK)\ \cup\ ([B]+3^LK).
\]
The two pieces have different first ternary digits. In particular,
the two-word code is prefix-free, its infinite addresses are
unique, and these similarities have disjoint images on $K$.
Let $d>0$ solve $3^{-md}+3^{-Ld}=1$. Existence and uniqueness
follow from strict monotonicity: the left side decreases from
$2$ to $0$ as $d$ increases from $0$ to infinity.

Give the independent choices $A,B$ probabilities $3^{-md}$ and
$3^{-Ld}$. The induced measure $\mu$ on $K$ assigns mass
$3^{-ds}$ to an address cylinder whose completed blocks have
total ternary length $s$. Consider a $3$-adic ball of radius
$3^{-n}$ that meets $K$, so its first $n$ ternary digits are
fixed. The distinct first digits of $A$ and $B$ allow one to
parse from these digits every completed block before depth $n$.
That completed-block prefix is shared by all addresses in the
ball. Its length $s$ satisfies $s\ge n-L$, since each block
has length at most $L$. Hence
\begin{equation}\label{familypos:eq:mass}
                \mu(B)\le3^{-ds}\le3^{dL}3^{-dn}.
\end{equation}
The mass distribution principle yields $\dim_H K\ge d$.

For the reverse inequality, stop each address at the first
completed block of total length at least $n$. The stopped
lengths lie in $[n,n+L-1]$. These finitely many address cylinders
partition the address space, so the sum of their masses is $1$.
Their image diameters are at most $3^{-s}$, and thus the sum
of the $d$th powers of those diameters is at most $1$.
They cover $K$ at arbitrarily small scales. It follows that
$\mathcal H^d(K)\le1$ and $\dim_H K\le d$.

Appending the reachable entry prefix maps $K$ by a fixed
$3$-adic similarity into $C(1,M_k)$, and therefore preserves
its dimension. We have proved the subsystem assertion in
Theorem~\ref{familypos:main}; Theorem~\ref{thm:family} supplies
$M_k\notin\mathcal V$.

Finally $L_k\le4m_k$, with equality only for $k=0$.
Since $\vartheta^{-1}+\vartheta^{-4}=1$, at
$d=(\log_3\vartheta)/m_k$ the left side of
\eqref{familypos:eq:root} is at least $1$, and is strictly
larger for $k\ge1$. Monotonicity proves
\eqref{familypos:eq:theta} and its asserted strictness.

\begin{remark}[Scope]
This result describes an explicit positive-dimensional subset for
every member of the sporadic nonquotient family. Its lower bound
depends on $k$; no uniform positive lower bound or exact dimension
of the full intersection is asserted. It neither determines the
sharp asymptotics of $\MC(X)$ nor excludes the point $1$ from
the exceptional set for powers of two.
\end{remark}
